\documentclass[conference]{IEEEtran}
\IEEEoverridecommandlockouts
\usepackage{cite}
\usepackage{amsmath,amssymb,amsfonts}
\usepackage{algorithmic}
\usepackage{graphicx}
\usepackage{textcomp}
\usepackage{xcolor}
\usepackage{hyperref}
\usepackage{orcidlink}
\usepackage[utf8]{inputenc}
\def\BibTeX{{\rm B\kern-.05em{\sc i\kern-.025em b}\kern-.08em
    T\kern-.1667em\lower.7ex\hbox{E}\kern-.125emX}}
\begin{document}

\title{Conversión de Señales mediante Circuitos con Amplificadores Operacionales}\\

\author{%
    \IEEEauthorblockN{Ana Freitez\, \orcidlink{0009-0009-6506-499X}}
    \IEEEauthorblockA{\textit{Facultad de Ingeniería} \\ 
    \textit{Universidad Latina de Panamá} \\ 
    Ciudad de Panamá, Panamá}
    \and
    \IEEEauthorblockN{Cristell Cedeño\, \orcidlink{0009-0002-2617-663X}}
    \IEEEauthorblockA{\textit{Facultad de Ingeniería} \\ 
    \textit{Universidad Latina de Panamá} \\ 
    Ciudad de Panamá, Panamá}
}

\maketitle

\begin{abstract}
% Resumen breve (150-250 palabras): objetivo del proyecto, circuito implementado
% (integrador y derivador en cascada con amplificadores operacionales),
% metodología utilizada (simulación y/o implementación física) y principales
% resultados obtenidos.
[Completar resumen del proyecto aquí.]
\end{abstract}

\begin{IEEEkeywords}
amplificador operacional, integrador, derivador, procesamiento de señales, circuitos analógicos
\end{IEEEkeywords}

\section{Introducción}
% Contexto general del tema, motivación del proyecto y objetivo principal.
% Mencionar brevemente qué problema resuelve un circuito integrador/derivador
% en cascada y por qué es relevante estudiarlo.
[Completar introducción.]

\subsection{Objetivos}
\begin{itemize}
    \item Objetivo general: [Completar]
    \item Objetivos específicos:
    \begin{itemize}
        \item [Completar]
        \item [Completar]
    \end{itemize}
\end{itemize}

\section{Marco Teórico}
Un amplificador operacional es un circuito integrado muy versátil que permite diseñar múltiples circuitos electrónicos: desde comparadores de voltaje y amplificadores de señal, hasta filtros, sumadores y diferenciadores. Se caracteriza por tener dos entradas y una salida, las entradas del amplificador están polarizadas, lo que quiere decir que tendremos:
\begin{itemize}
	\item{Una entrada inversora (-).}
	\item{Una entrada no-inversora (+).}
	\item{Y una salida que está dada por la diferencia de las dos entradas multiplicadas por un factor.}
\end{itemize}

También se caracterizan por ser muy compactos, con un nivel de ganancia muy alto, por lo tanto, \textbf{el nivel de ganancia es el factor por el cual se multiplica la diferencia de las dos entradas.}

\begin{equation}
	V_{out} = A_{OL}(V_+, - V_-)
	\label{eq:amplificador}
\end{equation}


\subsection{Amplificador Integrador}

El amplificador integrador es una configuración especial de los amplificadores operacionales. Este circuito entrega una señal de salida, la cual es la integral en el tiempo de la señal de entrada.

Su principio de funcionamiento se basa en el amplificador operacional inversor (a excepción de la resistencia de retroalimentación) y en la operación matemática de integración.

\textbf{Fórmula voltaje de salida del amplificador integrador}

\begin{equation}
V_{out}(t) = -\frac{1}{RC}\int V_{in}(t)\,dt
\label{eq:integrador}
\end{equation}

\subsection{Amplificador Derivador}
El amplificador derivador es una configuración especial de los amplificadores operacionales. Este circuito tiene la capacidad de realizar la función matemática de derivación a una señal de entrada.

El principio de su funcionamiento de este amplificador se basa en la función matemática de la derivada, por ejemplo:

\begin{itemize}
	\item [Si introducimos una señal senoidal a la entrada del amplificador, este procesa esa señal y entrega en la salida una señal cosenoidal con una magnitud positiva.]
	
	\item[Si ingresamos en la entrada una señal triangular, en la salida obtendremos una señal cuadrada, ya que es la derivada de dicha señal triangular.]
\end{itemize}

\textbf{Fórmula de salida del amplificador derivador ideal}
\begin{equation}
V_{out}(t) = -RC\,\frac{dV_{in}(t)}{dt}
\label{eq:derivador}
\end{equation}

\subsection{Consideraciones de Diseño Práctico}

% Mencionar por qué se añaden componentes adicionales (resistencia en paralelo
% al capacitor del integrador, resistencia en serie y capacitor en paralelo
% en el derivador) para evitar saturación y ruido en la implementación real.
[Completar.]

\section{Metodología}

\subsection{Diseño del Circuito}
% Describir el diseño propuesto: etapa integradora seguida de etapa
% derivadora en cascada. Referenciar la Fig.~\ref{fig:diagrama}.
[Completar descripción del diseño.]

\begin{figure}[htbp]
\centerline{\includegraphics[width=0.9\linewidth]{diagrama_circuito.png}}
\caption{Diagrama esquemático del circuito integrador-derivador en cascada.}
\label{fig:diagrama}
\end{figure}

\subsection{Materiales y Equipo}
\begin{table}[htbp]
\caption{Lista de Materiales y Equipos}
\label{tab:materiales}
\centering
\begin{tabular}{l l c}
\hline
\textbf{Componente} & \textbf{Especificación} & \textbf{Cantidad} \\
\hline
Amplificador operacional & 741 & 2 \\
Resistencia & 10 k$\Omega$ & 3 \\
Resistencia & 100 k$\Omega$ & 1 \\
Capacitor & 0.1 $\mu$F & 3 \\
Fuente de alimentación & Bateria 9V & 1 \\
Protoboard & -- & 1 \\
Generador de señales & -- & 1 \\
Osciloscopio & -- & 1 \\
\hline
\end{tabular}
\end{table}

\subsection{Procedimiento}
% Describir paso a paso cómo se realizó la simulación y/o el montaje físico.
\begin{enumerate}
    \item [Completar]
    \item [Completar]
    \item [Completar]
\end{enumerate}

\section{Resultados}

\subsection{Simulación}
% Insertar capturas del circuito simulado y del osciloscopio virtual.
[Completar descripción de los resultados de simulación.]

\begin{figure}
	\centering
	\includegraphics[width=0.9\linewidth]{../Esquematico}
	\caption{Esquemático simulado en Proteus 8 Professional}
	\label{fig:sim_esq}
\end{figure}
\begin{figure}
	\centering
	\includegraphics[width=0.7\linewidth]{../Simulacion_osciloscopio}
	\caption{Formas de onda obtenidas en el osciloscopio virtual: señal de entrada, salida del integrador y salida del derivador.}
	\label{fig:sim_osc}
\end{figure}

\subsection{Implementación Física}

\begin{figure}
	\centering
	\includegraphics[width=0.7\linewidth]{../Montaje_fisico}
	\caption{Montaje físico del circuito en protoboard}
	\label{fig:montajefisico}
\end{figure}

\begin{figure}
	\centering
	\includegraphics[width=0.7\linewidth]{../Proyecto_Circuitos2/PNG/Captura_osciloscopio_real}
	\caption{Captura del osciloscopio real: a) señal de salida del integrador b) señal de salida del derivador}
	\label{fig:osc_real}
\end{figure}

\subsection{Datos Medidos}
% Completar la tabla con los valores obtenidos en simulación y/o en el
% laboratorio para distintas frecuencias de entrada.
\begin{table}[htbp]
\caption{Amplitudes de Salida Medidas en Función de la Frecuencia de Entrada}
\label{tab:datos}
\centering
\begin{tabular}{c c c c}
\hline
\textbf{$f$ (Hz)} & \textbf{$V_{in}$ (V)} &
\textbf{$V_{out,int}$ (V)} & \textbf{$V_{out,der}$ (V)}\\
\hline
25    & 4 & 4.400 V & 444.0 mV \\
50    & 4 & [Completar] & [Completar]\\
100   & 4 & [Completar] & [Completar] \\
200   & 4 & [Completar] & [Completar] \\
500   & 4 & [Completar] & [Completar] \\
1000  & 4 & [Completar] & [Completar] \\
\hline
\end{tabular}
\end{table}

\section{Análisis de Resultados}
% Comparar los resultados de simulación con los resultados físicos.
% Explicar cualquier diferencia observada (ruido, saturación, tolerancias
% de componentes, etc.) y relacionarlos con la teoría de la Sección II.
[Completar análisis comparando teoría, simulación e implementación física.]

\section{Conclusiones}
% Resumir los hallazgos principales y si se cumplieron los objetivos
% planteados en la introducción.
[Completar conclusiones.]

\begin{thebibliography}{00}
\bibitem{b1} [https://amplificadores.info/amp-op/derivador]
\bibitem{b2} [https://amplificadores.info/amp-op/integrador]
\bibitem{b3} [https://amplificadores.info/amp-op]
\bibitem{b3} [Completar referencia 3.]
\end{thebibliography}

\end{document}